\documentclass{article}

\usepackage{amsmath,amssymb}
\usepackage{xurl}
\usepackage[hidelinks]{hyperref}
\setlength{\emergencystretch}{2em}

\input{command}

\title{The Phase Index in the Weighted Ces\`aro Formula, Wolf Equation~(6.15)}
\author{}
\date{2026-08-16}

\begin{document}
\maketitle
\gapnote{local-correction}{resolved}

\section{The source sentence}

Let $T:\MN{D}\to\MN{D}$ be a trace-preserving positive linear map. After
expressing the asymptotic map $T_\infty$ as a Ces\`aro mean in
Equation~(6.14), Wolf writes
\begin{align}
    T_\phi
        &=\lim_{N\to\infty}\frac{1}{N}
          \sum_{n=1}^{N}
          \sum_{k:\lvert\lambda_k\rvert=1}
          (\overline{\lambda_k}T)^n.
    \label{eq:wolf_ch6_eq615_deduplicated_phases_block_indexed_formula}
\end{align}
This is Equation~(6.15) of Ref.~\cite{Wolf2012Quantum}. The local source is
\path{Notes/WolfNoteTexSource/ch06_spectral_properties.tex},
lines~258--264.

In this chapter, $k$ indexes the Jordan blocks of the superoperator
$\widehat T$. Equations~(6.4)--(6.5) of Ref.~\cite{Wolf2012Quantum}, at
lines~111--139 of the local source, use
\begin{align}
    \widehat T
        &=X\left(\bigoplus_{k=1}^{K}J_k(\lambda_k)\right)X^{-1}.
    \label{eq:wolf_ch6_eq615_deduplicated_phases_jordan_decomposition}
\end{align}
The number of blocks carrying an eigenvalue $\lambda$ is its geometric
multiplicity, so distinct block indices may carry the same eigenvalue
\cite[Chapter~6]{Wolf2012Quantum}.

\section{Correction forced by multiplicities}

The summand $(\overline{\lambda_k}T)^n$ depends only on $\lambda_k$, not on
the block indexed by $k$. Consequently, the literal block-indexed sum in
(\ref{eq:wolf_ch6_eq615_deduplicated_phases_block_indexed_formula}) counts
each peripheral eigenvalue with its geometric multiplicity.

For a counterexample, take the identity channel $T=\one$ on $\MN{D}$ with
$D>1$. Every matrix is an eigenvector with eigenvalue $1$, and the Jordan
decomposition of $\widehat T$ has $D^2$ one-dimensional blocks carrying that
eigenvalue. The literal right-hand side becomes
\begin{align}
    \lim_{N\to\infty}\frac{1}{N}
        \sum_{n=1}^{N}D^2\one
        &=D^2\one,
    \label{eq:wolf_ch6_eq615_deduplicated_phases_identity_block_sum}\\
    D^2\one
        &\neq \one=T_\phi.
    \label{eq:wolf_ch6_eq615_deduplicated_phases_identity_contradiction}
\end{align}
Thus Equation~(6.15) of Ref.~\cite{Wolf2012Quantum} is false when $k$ is
read literally as a Jordan-block index.

Let $\mathcal P(T)$ denote the set of distinct eigenvalues of $T$ on the unit
circle. The corrected formula is
\begin{align}
    T_\phi
        &=\lim_{N\to\infty}\frac{1}{N}
          \sum_{n=1}^{N}
          \sum_{\lambda\in\mathcal P(T)}
          (\overline{\lambda}T)^n.
    \label{eq:wolf_ch6_eq615_deduplicated_phases_distinct_phase_formula}
\end{align}
Each peripheral phase is counted once.

This reading agrees with the surrounding spectral decomposition. Peripheral
Jordan blocks of a trace-preserving positive map are one-dimensional
\cite[Chapter~6]{Wolf2012Quantum}. Hence the peripheral subspace is the
direct sum of the eigenspaces associated with the distinct peripheral
eigenvalues. If $Y$ belongs to the eigenspace of $\lambda$, then the summand
indexed by $\mu\in\mathcal P(T)$ satisfies
\begin{align}
    (\overline{\mu}T)^n(Y)
        &=(\overline{\mu}\lambda)^nY.
    \label{eq:wolf_ch6_eq615_deduplicated_phases_action_on_eigenspace}
\end{align}
Its Ces\`aro mean obeys
\begin{align}
    \lim_{N\to\infty}\frac{1}{N}
        \sum_{n=1}^{N}(\overline{\mu}\lambda)^n
        &=
        \begin{cases}
            1,&\mu=\lambda,\\
            0,&\mu\neq\lambda.
        \end{cases}
    \label{eq:wolf_ch6_eq615_deduplicated_phases_phase_orthogonality}
\end{align}
Therefore one summand for each distinct peripheral eigenvalue recovers the
peripheral projection.

The neighboring expressions for $\widehat T_\infty$, $\widehat T_\phi$, and
$\widehat T_\varphi$ in Equations~(6.11)--(6.13) of
Ref.~\cite{Wolf2012Quantum} are also indexed by Jordan blocks, but their
summands contain the block projections $P_k$. Those formulas retain block
dependence and are correct as printed. Only Equation~(6.15), where the
summand no longer depends on the block, requires a deduplicated index.

\section{Formalization}

The formal development averages over the set of eigenvalues of modulus one,
so each distinct phase occurs once.\footnote{Formal declarations:
    \leanid{Module.End.weightedCesaroMean} for the average, and
    \leanid{IsPositiveMap.tendsto_weightedCesaroMean_peripheralProjection},
    \leanid{IsPositiveMap.tendsto_endEquiv_weightedCesaroMean_peripheralProjection},
    \leanid{IsPositiveMap.tendsto_weightedCesaroMean_toFinset_peripheralProjection},
    and
    \leanid{IsPositiveMap.tendsto_endEquiv_weightedCesaroMean_toFinset_peripheralProjection}
    in \doclink{QICLean/Channel/Peripheral/WeightedCesaro}.}
This correction changes no hypothesis or intended conclusion of
Equation~(6.15) of Ref.~\cite{Wolf2012Quantum}; it changes only the
multiplicity convention for the inner index.

\section{Conclusion}

Equation~(6.15) of Ref.~\cite{Wolf2012Quantum} requires a local correction:
the inner sum must range over distinct peripheral eigenvalues, not peripheral
Jordan blocks. The identity channel gives the contradiction
(\ref{eq:wolf_ch6_eq615_deduplicated_phases_identity_contradiction}) under
the literal block-indexed reading. The deduplicated formula
(\ref{eq:wolf_ch6_eq615_deduplicated_phases_distinct_phase_formula}) has the
source's intended hypotheses and conclusion.

\bibliographystyle{alpha}
\bibliography{references}

\end{document}